\section{Atención a la diversidad}\label{sec:diversidad}
El decreto regula las medidas de apoyo y las condiciones de accesibilidad y diseño universal, y se rige por los principios de normalización e inclusión \parencite[art.~32, pp.~27--28]{decreto64}. La propuesta sigue las pautas del Diseño Universal para el Aprendizaje \parencite{cast2024}.

\begin{tablaSAlarga}{colspec={Q[l,wd=100pt]X[l]}, row{1}={bg=urjc, fg=white, font=\footnotesize\bfseries}, column{1}={font=\footnotesize\bfseries}}
  Principio DUA & Medidas para todo el grupo \\
  Implicación & Reto real con destinatario; elección del objeto a medir; roles rotatorios; error tratado como oportunidad. \\
  Representación & Explicación magistral, construcción dinámica y material escrito; glosario de geometría disponible durante toda la secuencia; ejemplos con el propio centro. \\
  Acción y expresión & Productos variados (libro, exposición, escrito, mapa conceptual); clinómetro casero o digital; plantillas de planificación y Kanban. \\
\end{tablaSAlarga}

\begin{tablaSAlarga}{colspec={Q[l,wd=100pt]X[2,l]X[1,l]}, row{1}={bg=urjc, fg=white, font=\footnotesize\bfseries}, column{1}={font=\footnotesize\bfseries}}
  Necesidad & Medidas y actividades & Elementos afectados \\
  Refuerzo (según Q1--Q4) & En S13, ejercicios resueltos paso a paso y parejas de apoyo sobre el criterio que indiquen los cuestionarios, seguidos de un segundo intento antes de S18. & Metodología; criterios 1.2, 2.1, 5.2 y 7.1, sin modificarlos \\
  Apoyo en la ficha F4 & Tarjetas de ayuda que el equipo pide cuando las necesita, sin separar a nadie de su grupo. & Metodología; criterio 7.2, sin modificarlo \\
  Ampliación y alto rendimiento & En S13, panel inclinado óptimo o más de un obstáculo en el modelo. Con los datos del socio, estimar el radio de la Tierra. & Contenidos de profundización; criterios 3.1, 6.1 \\
  TDAH & Tareas en bloques breves con hitos en Kanban, rotación ordenada para que tenga el rol de medición en la salida al patio (letra D) y más tiempo si lo necesita en el examen. & Tiempos y organización; adecuación del examen \\
  Dislexia & Enunciados en tipografía clara, lectura en voz alta, sin penalizar la ortografía en el examen y exposición oral como alternativa para la reflexión. El rol de comunicación le corresponde en el turno~1 (letra D), con textos breves y corrector. & Instrumentos y formato \\
  Incorporación reciente & Revisión de los conocimientos previos que indique el diagnóstico y compañero de referencia en el equipo. & Apoyo inicial \\
\end{tablaSAlarga}

En ningún caso se rebajan los criterios de evaluación: lo que cambian son los apoyos, los tiempos y los formatos, conforme al art.~32 del decreto.
